\documentclass[./main.tex]{subfiles}
\begin{document}

% Slide # 4
\begin{frame}[t, label=slide04]
        % Title
        \frametitle{Tipping mechanisms}

        % Body
        \footnotesize

        We want to differentiate between bifurcation-induced (B-tipping) from rate-induced (R-tipping) critical transitions.

        \vspace{0.5cm}

        The former type (usually) happens when the parameter shift is too abrupt. 

        \vspace{0.5cm}

        Suppose $\Lambda_{\varepsilon}(\mu_{-},\mu_{+}) = \frac{\mu_{-}+\mu_{+}}{2} + \frac{\mu_{+} - \mu_{-}}{2}\tanh(\varepsilon(t-t_{0}))$:

        \begin{columns}[T]
                
                \column{0.5\textwidth}

                \begin{figure}[H]
                        \centering 
                        \includegraphics<2-3>[keepaspectratio, width=\textwidth]{./figures/slide_04_1.png}%
                \end{figure}

                \hfill

                \column{0.5\textwidth}

                \begin{figure}[H]
                        \centering 
                        \includegraphics<3>[keepaspectratio, width=\textwidth]{./figures/slide_04_2.png}%
                \end{figure}

        \end{columns}

\end{frame}

\end{document}
